\section{A reference implementation of CONSTRUCT-MP}
\label{app:code}

This appendix gives a Python sketch of CONSTRUCT-MP applied to a toy operator algebra: bit-flip-symmetric programs over $\mathbb{F}_2^n$. The complete code (approximately 80 lines including comments and demonstration), along with the reactor-physics and equivariant-ML instantiations, will be released on Zenodo at acceptance.

\begin{lstlisting}[style=pythonstyle, caption={Reference implementation of CONSTRUCT-MP for a toy bit-flip algebra.}]
"""
NOETHER reference implementation: CONSTRUCT-MP on a toy bit-flip algebra.
Demonstrates the four-step procedure of Section 4.2.
"""
from dataclasses import dataclass, field
from typing import Callable, Dict, List, Set, Tuple

@dataclass(frozen=True)
class Operator:
    name: str
    block: str   # one of {"G","O_le","T*","T_rev*","L*","D*","E*"}

@dataclass(frozen=True)
class Invariant:
    operator: Operator
    description: str

@dataclass(frozen=True)
class MR:
    invariant: Invariant
    template: str

@dataclass
class MetaPattern:
    block: str
    members: Set[MR] = field(default_factory=set)

CANONICAL_ORDER = ["G","O_le","T*","T_rev*","L*","D*","E*"]   # Def. 3

def extract_invariants(generators: List[Operator]) -> Dict[str, List[Invariant]]:
    inv_by_block = {b: [] for b in CANONICAL_ORDER}
    for g in generators:
        if g.block == "G":
            inv_by_block["G"].append(
                Invariant(g, f"P(g.x) = rho(g).P(x) for g={g.name}"))
        elif g.block == "O_le":
            inv_by_block["O_le"].append(
                Invariant(g, f"x1 <=_theta x2 => P(x1) <=_Y P(x2) for theta={g.name}"))
    return inv_by_block

def translate(iota: Invariant) -> MR:
    return MR(iota, template=iota.description)

def construct_mp(generators: List[Operator]) -> List[MetaPattern]:
    inv_by_block = extract_invariants(generators)
    mps = []
    for block in CANONICAL_ORDER:
        mrs = {translate(i) for i in inv_by_block[block]}
        if mrs:
            mps.append(MetaPattern(block=block, members=mrs))
    return mps

if __name__ == "__main__":
    Z2 = Operator("bit_flip", "G")
    perm = Operator("permutation_S3", "G")
    scale_x = Operator("input_scaling", "O_le")

    mp_set = construct_mp([Z2, perm, scale_x])
    for mp in mp_set:
        print(f"MetaPattern in block {mp.block}:")
        for mr in mp.members:
            print(f"  - {mr.template}")
\end{lstlisting}

The reactor-physics and equivariant-ML instantiations follow the same structure, replacing the toy generators with the operators of $\mathcal{A}_{\mathrm{Boltz}}$ and $\mathcal{A}_{\mathrm{equi}}$ respectively.

%% =========================================================
\section{Construct-trace consistency check on hand-crafted block-targeted mutants}
\label{app:augmented-stratum}

This appendix reports a \emph{construct-trace consistency check} on
the five operative NOETHER blocks ($O_{\le}$,
$\mathcal{T}^{*}_{\mathrm{rev}}$, $\mathcal{D}^{*}$, $\mathcal{E}^{*}$,
$\mathcal{B}^{*}_{\mathrm{rel}}$) that PIT 1.7.4's stock mutator
catalogue does not exercise on the
\S\ref{subsec:test-design} substrate. The exercise is not an
independent fault-detection comparison: by construction, each mutant
violates a specific Set~N MR's targeted invariant, so a non-trivial
Set~N detection rate is implied by the framework's
\texttt{CONSTRUCT-MP} design rather than measured by an independent
oracle. We report the results to (i) verify that the per-mutant
\texttt{jir}/\texttt{jor} pipeline does in fact catch the intended
invariant violations end-to-end, and (ii) document the secondary
question of whether Set~G's GP-evolved MRs incidentally detect the
same mutants. The headline H3a verdicts in \S\ref{subsec:pooled-headtohead}
do not rely on the numbers in this appendix.

\paragraph{Mutant authoring.}
For each of the five uncovered operative blocks we hand-crafted $5$
mutants ($25$ total) against the \S\ref{subsec:test-design} SUTs whose
Set~N MR catalogue contains an MR targeting that block. Each mutant
is a complete \texttt{MathSignalClass.java} replacement with one
method body altered to violate the targeted block's canonical-form
invariant (per-mutant rationale in
\texttt{docs/block\_coverage\_augmentation\_rationale.md} of the
experiment repository; generator and runner in commits
\texttt{c2030d9} and \texttt{2712564}; codegen fixes to remove
forward-reference and identity-fallback regressions in commit
\texttt{f7edfa7}).\footnote{A residual codegen limitation, documented
in \texttt{docs/codegen\_fix\_audit.md} of the experiment repository
but not fixed in the current build, is that
\texttt{parse\_jir\_to\_transforms} does not handle the
integer-equality jir pattern \texttt{i\_n\_f == -i\_n\_s} used by the
\texttt{D\_reciprocal\_exp} MR on \texttt{powerSig}; under the
current codegen the MR's follow-up exponent falls back to identity
($n_{f} = n_{s}$) and the jor \texttt{|o\_f - 1/o\_s| < 1e-4}
degenerates into a coarse ``output is not its own reciprocal''
detector. The construct-trace kill verdict for the
\texttt{powerSig} mutants reported in
Table~\ref{tab:augmented-stratum} remains robust to removing the
$D\_reciprocal\_exp$ MR (every \texttt{powerSig} augmented mutant
is also killed by at least one of \texttt{E\_power\_of\_power},
\texttt{L\_scale\_base}, \texttt{T\_exp\_step}, or
\texttt{G\_odd\_negate}, per the per-MR audit in the experiment
repository), so the appendix's per-block kill flag is unaffected;
nonetheless the MR-level specificity question of which MR fires
on which mutant is partly clouded by this degeneracy, and a
codegen extension to handle integer-equality jir patterns is
committed as a separate follow-up.} Mutant evaluation uses \texttt{mvn test} on the
substituted SUT under each existing Set~N and Set~G MR test suite;
kill verdict is "killed by any Set~N MR" (Definition~A of
\texttt{docs/d1d2\_methodology.md} \S4 in the experiment repository).

\paragraph{Per-block results.}
Table~\ref{tab:augmented-stratum} reports the per-block kill counts.
Set~N detects $25/25$ of the hand-crafted mutants
(Wilson 95\% CI $[0.867,\,1.000]$); since each mutant was
constructed to violate the block targeted by a Set~N MR, this is the
\emph{design-implied} outcome and should be read as confirming the
construct-trace consistency of the pipeline rather than as a measure
of independent fault-detection. Set~G's
$12/25 = 0.480$ ($[0.300,\,0.665]$) is a more informative number: it
quantifies the incidental reach of GP-evolved generic MRs on the
same hand-crafted mutants, with $0$ Set-G-only kills. Two blocks
($\mathcal{T}^{*}_{\mathrm{rev}}$ and
$\mathcal{B}^{*}_{\mathrm{rel}}$, both targeting the
\texttt{isSequence} SUT) have Set~G $0/5$ by
absence-of-MR-coverage: the 30-min GenMorph rerun produced no
\texttt{isSequence} MR (see \S\ref{subsec:empirical-threats}~(d) on
GenMorph's structural N/A pattern).

\begin{table}[h]
\centering
\caption{Construct-trace consistency check on hand-crafted block-targeted
mutants ($n = 5$ per block, $25$ total). \textbf{Set~N's column is
design-implied}: each mutant is constructed against a specific Set~N
MR's targeted invariant, so Set~N's near-$100\%$ detection is a
pipeline-correctness check, not independent fault-detection evidence
(see \nameref{para:construct-trace-not-headtohead} below).
Set~G's $12/25 = 0.480$ measures the incidental reach of GP-evolved
generic MRs on the same mutants and is not design-implied.
Wilson 95\% intervals reported. Two blocks have Set~G $0/5$ from
GenMorph's structural absence of an \texttt{isSequence} MR at the
$30$-min budget. \textbf{This table is excluded from the H3a.1
evidence base (\S\ref{subsec:pooled-headtohead});} the H3a.1 verdict
rests on the pre-registered PIT-covered three-block substrate only.}
\label{tab:augmented-stratum}
\small
\begin{tabular}{lr l l r l r r r r}
\toprule
Block & $n$ & \emph{Set~N kills (CTT)} & \emph{rate (CTT)} & Set~G kills & rate (95\% CI) & both & $N$-only & $G$-only & neither \\
\midrule
$O_{\le}$                                  & 5 & \emph{5} & \emph{$1.000^{\,\dagger}$} & 3 & $0.600$ $[0.231,\,0.882]$ & 3 & 2 & 0 & 0 \\
$\mathcal{T}^{*}_{\mathrm{rev}}$           & 5 & \emph{5} & \emph{$1.000^{\,\dagger}$} & 0 & $0.000$ $[0.000,\,0.434]$ & 0 & 5 & 0 & 0 \\
$\mathcal{D}^{*}$                          & 5 & \emph{5} & \emph{$1.000^{\,\dagger}$} & 4 & $0.800$ $[0.376,\,0.964]$ & 4 & 1 & 0 & 0 \\
$\mathcal{E}^{*}$                          & 5 & \emph{5} & \emph{$1.000^{\,\dagger}$} & 5 & $1.000$ $[0.566,\,1.000]$ & 5 & 0 & 0 & 0 \\
$\mathcal{B}^{*}_{\mathrm{rel}}$           & 5 & \emph{5} & \emph{$1.000^{\,\dagger}$} & 0 & $0.000$ $[0.000,\,0.434]$ & 0 & 5 & 0 & 0 \\
\midrule
\textit{aggregate} & 25 & \emph{25} & \emph{$1.000^{\,\dagger}$} & 12 & $0.480$ $[0.300,\,0.665]$ & 12 & 13 & 0 & 0 \\
\bottomrule
\end{tabular}

\smallskip
\footnotesize
\textit{CTT = construct-trace test (design-implied).}
$^\dagger$ Set~N columns are italicised to mark that the rate is design-implied by mutant authoring (each mutant constructed against a specific Set~N MR's targeted invariant), not an independent fault-detection measurement. Do not compare $1.000^{\,\dagger}$ against Set~G's $0.480$ as if they were commensurable; the $1.000$ value is a pipeline-correctness check, the $0.480$ value is Set~G's incidental reach on the same mutants. Comparative H3a.1 evidence comes from the PIT-covered substrate (Table~\ref{tab:per-block-headtohead}, \S\ref{subsec:pooled-headtohead}).
\end{table}

\paragraph{Multiple-comparison correction.}
A Bonferroni-corrected Fisher exact test at the per-block
significance level $\alpha / 5 = 0.01$ would be appropriate if these
were independent measurements of Set~N's reach. Since the test is
design-implied and the per-block sample size $n = 5$ is below the
threshold at which any cell of the $2 \times 2$ table can reach
$p < 0.01$ in either direction (Fisher exact lower bound under
$n_{1} = n_{2} = 5$ with all-or-nothing splits is $p \approx 0.008$,
attainable only at the
$0/5$-vs-$5/5$ extreme; both
$\mathcal{T}^{*}_{\mathrm{rev}}$ and
$\mathcal{B}^{*}_{\mathrm{rel}}$ rows reach this extreme but only
via Set~G's absence-of-coverage), the construct-trace reading is
the only defensible interpretation.

\paragraph{Why this is not a head-to-head test of H3a.1.}
\label{para:construct-trace-not-headtohead}
H3a.1 (Set~N's per-block detection on D1) is a question about
\emph{generic} algebra-disrupting mutations on the substrate, of
which PIT 1.7.4 produces a subset on three operative blocks (G,
$\mathcal{L}^{*}$, $\mathcal{T}^{*}$). The hand-crafted mutants of
Table~\ref{tab:augmented-stratum} are not drawn from a generic
mutation distribution: they were specifically authored to violate
the targeted invariant of a known Set~N MR. To use them as
evidence for H3a.1 on the five additional operative blocks would
constitute construct-trace circular validation. The
construct-trace check is reported here for transparency and to
verify the pipeline's end-to-end correctness on the previously
uncovered blocks; the H3a.1 verdict reported in
\S\ref{subsec:pooled-headtohead} rests on the pre-registered
PIT-covered three-block substrate only.

%% =========================================================
